% Einheit 2 von 5 – Kürzen und Erweitern (bank/brueche-dezimalzahlen/e2.jsonl, zusammenbau v0.1)
%% TODO Zweigzeile Teil 1: was hier gelernt wird (Satz in Schülersprache aus dem Einheitstitel) – Einheit 2
\einheitenkopf[e2]{Einheit 2 von 5 · Kürzen und Erweitern}
\zweigzeile{ab Kl. 5, je nach Buch bis Kl. 6 · P10 oft}
\setcounter{aufgabe}{13}

%% TODO Titel: Ich-kann-Satz fehlt in der Bank (Platzhalter: Kettenname) – Nr. 14
%% TODO Anweisung: ein Satz über der ersten Teilaufgabe fehlt in der Bank – Nr. 14
\begin{aufgabe}{Kürzen und Erweitern}
%% TODO \rechenplatz{n}: Schrittzahl des Grundfalls fehlt in der Bank (nur bei fester Schreibform, 2.3 b)
\begin{teile}
\teil $\frac{2}{5} \to \frac{6}{15}$ – mit welcher Zahl erweitert?
\teil Erweitere $\frac{3}{5}$ mit $4$.
\teil Kürze $\frac{14}{21}$ mit $7$.
\teil Kürze $\frac{16}{24}$ vollständig.
\teil Schreibe $\frac{3}{4}$ mit dem Nenner $20$.
\teil Mache gleichnamig: $\frac{1}{3}$ und $\frac{5}{9}$.
\teil Mache gleichnamig: $\frac{3}{4}$ und $\frac{2}{5}$.
\teil Schreibe $\frac{13}{5}$ als gemischte Zahl.
\teil Welcher Anteil der Figur ist grau? Gib ihn als vollständig gekürzten Bruch und in Prozent an. \hfill (P10 2014 OS) \\

\bruchrechteck{12}{16}
Bruch: \leerfeld Prozent: \leerfeld[\%]
\end{teile}
\end{aufgabe}

%% TODO Titel: Ich-kann-Satz fehlt in der Bank (Platzhalter: Kettenname) – Nr. 15
%% TODO Anweisung: ein Satz über der ersten Teilaufgabe fehlt in der Bank – Nr. 15
\begin{aufgabe}{gleichwertige Brüche am Streifen finden}
\begin{teile}
\teil Der obere Streifen zeigt $\frac{1}{2}$. Wie viele Teile musst du im unteren Streifen färben, damit gleich viel gefärbt ist?

\bruchrechteck{1}{2} \\ \bruchrechteck{0}{8}
\end{teile}
\end{aufgabe}

%% TODO Titel: Ich-kann-Satz fehlt in der Bank (Platzhalter: Kettenname) – Nr. 16
%% TODO Anweisung: ein Satz über der ersten Teilaufgabe fehlt in der Bank – Nr. 16
\begin{aufgabe}{Bruch als Geteilt-Aufgabe}
\begin{teile}
\teil Schreibe $7 : 9$ als Bruch.
\end{teile}
\end{aufgabe}

%% TODO Titel: Ich-kann-Satz fehlt in der Bank (Platzhalter: Kettenname) – Nr. 17
%% TODO Anweisung: ein Satz über der ersten Teilaufgabe fehlt in der Bank – Nr. 17
\begin{aufgabe}{Kürzen und Erweitern – Fehler finden, Begründen, Darstellungswechsel, Anwendung}
\begin{teile}
\teil Mia erweitert $\frac{2}{7}$ mit $3$ und schreibt $\frac{6}{7}$. Finde den Fehler.
\teil Erweitert man $\frac{1}{2}$ mit $3$, stehen größere Zahlen da. Ist der Bruch größer geworden? Begründe.
\teil Jedes Teil des Streifens wird in $3$ gleiche Teile geteilt. Welcher Bruch steht dann für den gefärbten Anteil?

\bruchrechteck{2}{5}
\leerfeld
\teil Eine Busfahrt dauert $45$ Minuten. Welcher Bruchteil einer Stunde ist das? Gib ihn gekürzt an. \leerfeld Stunde
\end{teile}
\end{aufgabe}
